% ============================================================
% DADOS GERAIS DA SINDICÂNCIA
% Arquivo: dados/dados-sindicancia.tex
%
% Este arquivo deve existir sempre.
% Ele concentra SOMENTE os dados gerais reutilizados por mais
% de um modelo/anexo.
%
% Regra do projeto:
% - Dados comuns ficam aqui.
% - Dados exclusivos de uma peça ficam em dados-anexo-*.tex.
% - Evite redefinir estas macros nos arquivos de anexos, pois
%   o LaTeX usa sempre o último valor carregado no main.tex.
% ============================================================

% ------------------------------------------------------------
% Data/local de abertura dos trabalhos
% Usada também como fallback quando um documento não informar data própria.
% ------------------------------------------------------------
\dia{00}
\mes{mês}
\ano{2026}
\cidade{CIDADE/UF}

% ------------------------------------------------------------
% Organização militar
% ------------------------------------------------------------
\om{ORGANIZAÇÃO MILITAR}
\omSuperior{ESCALÃO SUPERIOR}
\omComplemento{LINHA COMPLEMENTAR DA OM}
\omNomeHistorico{NOME HISTÓRICO DA OM}

% ------------------------------------------------------------
% Processo / portaria / objeto
% ------------------------------------------------------------
\nup{00000.000000/2026-00}

\portaria{Portaria nº 000, de 00 de mês de 2026, do Comandante da ORGANIZAÇÃO MILITAR}

\objeto{apurar os fatos constantes da portaria instauradora e seus anexos}

% ------------------------------------------------------------
% Autoridade instauradora / nomeante
% ------------------------------------------------------------
\autoridadeNome{AUTORIDADE INSTAURADORA}
\autoridadePosto{POSTO}
\autoridadeFuncao{Comandante da ORGANIZAÇÃO MILITAR}

% ------------------------------------------------------------
% Sindicante
% ------------------------------------------------------------
\sindicanteNome{SINDICANTE}
\sindicantePosto{POSTO/GRADUAÇÃO}
\sindicanteOM{\getOm}

% ------------------------------------------------------------
% Sindicado
% Se ainda não houver sindicado, deixe em branco ou comente,
% desde que o modelo usado permita ausência.
% ------------------------------------------------------------
\sindicadoNome{SINDICADO}
\sindicadoPosto{POSTO/GRADUAÇÃO}
\sindicadoSU{SUBUNIDADE/SEÇÃO}
\sindicadoOM{\getOm}

% Dados complementares do sindicado, usados principalmente no Anexo Q
% e também por outros modelos que precisem qualificar o sindicado.
\sindicadoIdentidade{000000000 ÓRGÃO/UF}
\sindicadoCPF{000.000.000-00}
\sindicadoNascimento{00 de mês de 0000}
\sindicadoNaturalidade{CIDADE/UF}
\sindicadoEstadoCivil{ESTADO CIVIL}
\sindicadoFiliacao{NOME DO PAI e NOME DA MÃE}
\sindicadoResidencia{ENDEREÇO COMPLETO}
\sindicadoCelular{(00) 00000-0000}

% ------------------------------------------------------------
% Escrivão
% Usado pelos Anexos F e G, quando houver designação de escrivão.
% ------------------------------------------------------------
\escrivaoNome{ESCRIVÃO}
\escrivaoPosto{POSTO/GRADUAÇÃO}

% ------------------------------------------------------------
% Campos genéricos opcionais
% ------------------------------------------------------------
\prazo{3 (três) dias úteis}
